% !PRINTED 720
\weberpage{720}

%<<CONT
whereby $D$ is defined as a positive number different from zero.
As regards the sums
\[
X = \overset{n}{\Sigma}\, \frac{\chi(n)}{n}
\]
it is to be recalled that $\chi(n)$ [by \S.\,193, (8)] is defined by
\weq{7}{n \equiv (-1)^\alpha 5^\beta \pmod{m}, \qquad
\chi(n) = (\pm 1)^\alpha \Theta^\beta,}
in which $\Theta$ denotes a $\nu$th root of unity.

But among the $\nu$th roots of unity the $\nu'$th roots of unity are
also contained, and among the characters $\chi$ also the characters for
the field $\Omega_m'$. We accordingly distinguish \gesperrt{primitive and
imprimitive characters of the field $\Omega_m$}, in that we designate as
primitive \gesperrt{those characters $\chi$ peculiar to the field $\Omega_m$,
i.\,e. those in which $\Theta$ is a primitive $\nu$th root of unity}.

Among the sums $X$ there occur also all the sums belonging to the field
$\Omega_m'$:
\[
X' = \overset{n}{\Sigma}\, \frac{\chi'(n)}{n}.
\]

If one now substitutes in (4) for $k, k'$ and $h_1, h_1'$ the
expressions formed in \S.\,195, (16), (17), and then cancels out the
sums $X'$, there results:
\weq{8}{\begin{aligned}
\pi^{\nu'} A &= 2 \cdot 2^{2^{\varkappa-4}(\varkappa-2)} \prod X_2 \\
D B &= 2^{2^{\varkappa-4}(\varkappa-2)} \prod X_1,
\end{aligned}}
\gesperrt{in which, however, the products $\prod$ now extend only to the sums
$X_1, X_2$ formed with the primitive characters $\chi_1, \chi_2$.}

In the expressions for the $X$ given in \S.\,194, (12):
\weq{9}{\begin{aligned}
X_1 &= -\frac{1}{2m} \sum_{1,\,m-1}^{s} f_1(r^s)
\log\left(\sin \frac{s\pi}{m}\right)^2 \\
X_2 &= -\frac{i\pi}{m^2} \sum_{1,\,m-1}^{s} s\, f_2(r^s)
\end{aligned}}
considerable simplifications now occur, under the assumption that $m$ is
a power of 2 and that $\chi_1, \chi_2$ are primitive characters.
